
\chapter{Superconductivity as a BEC of Pairs}
In chapter \ref{Ch_2} we have shown that superconductivity originates from the pairing of electrons into correlated states known as \textit{Cooper pairs}.  
Each pair carries charge \(2e\) and total spin zero\footnote{In the singlet case: two electrons with opposite spins $(\uparrow\downarrow)$ combine to give a total spin $J = 0$.} and thus behaves as a composite boson.  
When cooled below a critical temperature \(T_c\), these composite bosons condense into a single macroscopic quantum state characterized by long-range phase coherence.  
This condensed state constitutes the superconducting phase.

\section{Macroscopic Analogy}
From a macroscopic standpoint, superconductivity and Bose--Einstein condensation share the same fundamental feature: both arise from the occupation of a single quantum state by a macroscopic number of particles.  
The key distinction lies in the microscopic mechanism that forms the condensed particles:  

\begin{itemize}
    \item \textbf{Bose--Einstein condensate:} In a conventional Bose gas, the constituent particles are \emph{elementary bosons} (e.g.\ atoms or photons) that condense directly into the lowest energy state below a critical temperature.

    \item \textbf{Superconductor:} In a superconductor, the condensed entities are not elementary particles but \emph{composite bosons} (bound states of two electrons known as \emph{Cooper pairs}), each carrying charge \(2e\) and integer spin.
    The condensation of these composite bosons into a single coherent quantum state produces the superconducting phase.
\end{itemize}

Hence, superconductivity can be viewed as a \textit{Bose--Einstein condensation of Cooper pairs}.

The analogy is made even clearer through the order parameter formalism.  
In BEC, the macroscopic wavefunction of the condensate is expressed as 
\[
\Psi(\mathbf{r}) = \sqrt{n_0(\mathbf{r})}\, e^{i\phi(\mathbf{r})},
\]
where \(n_0\) is the condensate density and \(\phi\) its global phase.  
In a superconductor, the order parameter takes an analogous form,
\[
\Delta = |\Delta| e^{i\phi},
\]
whose magnitude \(|\Delta|\) measures the superconducting energy gap, and whose phase \(\phi\) defines the macroscopic coherence of the condensate.  
Below the critical temperature \(T_c\), \(\Delta\) acquires a finite value, signalling the onset of long-range phase order and the spontaneous breaking of the global \(\mathrm{U}(1)\) gauge symmetry.  
This broken symmetry gives rise to the collective quantum state responsible for superconductivity.

Hence, while the mechanisms of pairing and condensation differ, their consequence is the same: a \textit{phase-coherent macroscopic quantum state}.  \\
This correspondence naturally raises a question:
\begin{center}
\textit{Do the BCS and BEC descriptions represent two distinct phases, or merely two limits of a single underlying phenomenon?}
\end{center}

%\section{Weak- To Strong- Coupling Regime} %TO DO!! Choose one
\section{The BCS-BEC Crossover}

A unified description of superconductivity can be obtained by extending the Bose--Einstein condensation formalism to include systems of composite bosons whose dispersion relation depends on the strength of the underlying fermionic attraction. 

\subsection*{General expression for the critical temperature.}
The condensation temperature and the dimensionality dependence can be derived for a general dispersion relation of the form
\[
    \varepsilon_K = C_s K^s,
\]
where \(C_s\) is a dimensional constant and \(s\) defines how the excitation energy scales with momentum. For a $d$-dimensional system the critical temperature reads:
\begin{equation}
\boxed{
    T_c = \frac{C_s}{k_B}
    \left[
        \frac{s\, \Gamma(d/2)\, (2\pi)^d\, n_B}
             {2\, \pi^{d/2}\, \Gamma(d/s)\, g_{d/s}(1)}
    \right]^{s/d}}\, ,
    \label{eq:Tc_general_1}
\end{equation}
where \(n_B\) is the boson (or Cooper-pair) density and we recall $g_{d/s}$ is the Bose function of order $d/s$. Using the relation,
\[
g_{\sigma}(1) = \zeta(\sigma) \iff g_{d/s}(1) = \zeta(d/s) 
\]
\eqref{eq:Tc_general_1} can be equivalently written as:
\begin{equation}
\boxed{
    T_c = \frac{C_s}{k_B}
    \left[
        \frac{s\, \Gamma(d/2)\, (2\pi)^d\, n_B}
             {2\, \pi^{d/2}\, \Gamma(d/s)\, \zeta(d/s)}
    \right]^{s/d}}\, ,
    \label{eq:Tc_general_2}
\end{equation}

The condensation condition requires \(d/s > 1\), ensuring that the number integral converges [Eq.~\eqref{convergence}]. For the conventional case of quadratic dispersion, condensation is possible only if \(d > 2\), as we have already seen in section \ref{crit}. Hence, in strictly two-dimensional systems, the integral diverges and no finite critical temperature can exist.\\
Then, 
\begin{center}
    \textit{How can we observe superconductivity, or equivalently, Bose--Einstein condensation of pairs, in two-dimensional materials?}
\end{center}
The answer lies in the nature of the pair dispersion relation.  
In presence of a Fermi sea, the Pauli exclusion principle restricts the available momentum states for pairing, causing the centre-of-mass motion of the pairs to deviate from the free-particle quadratic law.  
As a result, the pair dispersion becomes approximately \emph{linear} rather than quadratic: this modification is directly linked to the strength of the underlying attractive interaction between fermions, which determines whether the system behaves in the weak- or strong-coupling limit.

\subsection*{Weak-coupling regime: BCS Limit}

In the weak-coupling limit, the attractive interaction between fermions is small compared to the Fermi energy, a condition quantified by the dimensionless ratio
\[
\frac{E_b}{\varepsilon_F} \ll 1,
\]
where \(E_b\) denotes the pair binding energy and \(\varepsilon_F\) the Fermi energy.  
In this regime, the binding energy of a Cooper pair is much smaller than the characteristic energy scale of the Fermi sea.  
The pairs are therefore highly delocalized in real space, with sizes much larger than the average interparticle distance, and they overlap extensively with one another.  
Because of this strong overlap, pair formation and condensation occur simultaneously at the same temperature \(T_c\), a defining feature of the \textit{BCS limit}:
\[
T_{\text{pair formation}} = 
T_{\text{condensation}} = T_c
\]

Microscopically, Cooper pairs form within the filled Fermi sea.  
However, the Pauli exclusion principle forbids the constituent fermions from occupying states already filled below the Fermi surface, restricting the available momentum states for pairing.  
This exclusion modifies the pair’s centre-of-mass motion: instead of moving freely as in vacuum, the pair interacts with the background Fermi sea.  
Consequently, its energy no longer follows the free-particle quadratic dependence on momentum but becomes approximately \textit{linear}:
\begin{equation}
    \varepsilon_K \approx \alpha(d)\, \hbar v_F\, K,
    \label{eq:linear_dispersion}
\end{equation}
where \(v_F\) is the Fermi velocity and \(\alpha(d)\) is a dimension-dependent coefficient, which values are \(\alpha(2) = 2/\pi\), \(\alpha(3) = 1/2\).  
This linear dispersion reflects the coupling of the pair’s motion to collective excitations near the Fermi surface and is a direct consequence of Pauli blocking.

Substituting \(s = 1\), $C_1 = \alpha(d)\, \hbar v_F$ into the generalized critical temperature formula~\eqref{eq:Tc_general_1}, we obtain another special case\footnote{For the full derivation, see Appendix \ref{app:BCS_BEC-derivation}, section \ref{app:Tc_linear_dispersion}.}:
\begin{equation}
\boxed{
    T_c = \frac{\alpha(d)\,\hbar v_F}{k_B}
    \left[
        \frac{\pi^{(d+1)/2}\, n_B}
             {\Gamma\!\left[(d+1)/2\right]\, g_d(1)}
    \right]^{1/d}}\,.
    \label{eq:Tc_linear_dispersion}
\end{equation}
This expression shows that for linear dispersion, the condensation temperature \(T_c\) remains finite for all \(d > 1\), thereby resolving the apparent paradox of superconductivity in quasi-two-dimensional systems.

\paragraph{Pair Breakup and the Finite Critical Temperature}
The previous analysis assumes the existence of ``unbreakable'' Cooper pairs, i.e.\ bound states that persist for all centre-of-mass momenta (CMM) from \(K = 0\) to \(K \to \infty\).  
In this idealized case, the integration over all momenta in the number equation leads to the result~\eqref{eq:Tc_linear_dispersion} for the condensation temperature.  
However, in a more realistic scenario, Cooper pairs can dissociate into two unpaired fermions once their total momentum exceeds a finite value \(K_0\).  
This critical momentum corresponds to the point where the pair binding energy vanishes,
\[
\Delta_{K_0} = 0,
\]
and defines the upper limit for the existence of stable pairs.

The breakup momentum \(K_0\) increases with the coupling strength:  
in the weak-coupling limit it is approximately
\[
K_0 \simeq \frac{\Delta_0}{a(d)\,\hbar v_F},
\]
and tends to infinity as the coupling becomes very strong (\(\Delta_0 \to \infty\)).  
Physically, \(K_0\) represents the maximum centre-of-mass momentum that a Cooper pair can sustain before dissociating into unpaired fermions.

By replacing the upper integration limit in the number equation with \(K_0\),  
the total boson number can be expressed in its general \(d\)-dimensional form as
\begin{equation}
N_{d\mathrm{D}} 
= \int_0^{\varepsilon_{K_0}} D_{d\mathrm{D}}(\varepsilon)\, f(\varepsilon)\, d\varepsilon
= \int_{|\mathbf K|\le K_0}\!\frac{d^d K}{(2\pi)^d}\, f(\varepsilon_K)
= \int_{|\mathbf K|\le K_0}\!\frac{d^d K}{(2\pi)^d}\,
      \frac{1}{e^{\beta(\varepsilon_K-\mu)} - 1}.
\end{equation}

From this modified number equation, it's possible to derive the asymptotic behaviour of the critical temperature in the limit of vanishing pair binding energy (\(\Delta_0 \to 0\)):
\begin{equation}
\boxed{
T_c \simeq
\frac{(d-1)\,[2\pi a(d)\hbar v_F]^d\, n_B}
{A_d\, k_B\, \Delta_0^{\,d-1}}
\quad \xrightarrow[\Delta_0 \to 0]{} \infty}\, ,
\label{eq:Tc_divergent}
\end{equation}
where \(A_2 = 2\pi\) and \(A_3 = 4\pi/3\).

Equation~\eqref{eq:Tc_divergent} shows that, as the pair binding energy \(\Delta_0\) approaches zero, the predicted critical temperature diverges.  
This unphysical result arises because, in this limit, all remaining pairs occupy the \(K = 0\) state at any temperature, implying that the system would be condensed at all \(T\).  
Physically, this corresponds to a situation in which the pair breakup threshold disappears, and every pair remains bound regardless of thermal excitation.

In practice, however, Cooper pairs do not remain bound for arbitrarily large centre-of-mass momenta.
Above a finite threshold $K_0$ the attractive interaction is no longer strong enough to keep the electrons paired, and the pair breaks apart into two unbound fermions.\\
When this physically unavoidable pair-breaking process happens, resulting in a more realistic \textit{binary boson–fermion mixture}, the artificial divergence of $T_c$ disappears.
The available phase space for bosonic excitations becomes finite, and the critical temperature is “regularized’’ to a finite value. Remarkably, this corrected $T_c$ turns out to be very close to the one predicted by the ideal Bose gas expression~\eqref{eq:Tc_linear_dispersion}, confirming that the divergence arose solely from the unphysical assumption of unbreakable pairs extending to infinite momentum.

\subsection*{Strong-coupling regime: BEC Limit}
In the strong-coupling limit, the attractive interaction between fermions dominates over the Fermi energy, a condition expressed by
\[
\frac{E_b}{\varepsilon_F} \gg 1.
\]
Here, the pair binding energy \(E_b\) is much larger than the characteristic energy scale of the Fermi sea.  
As a result, electrons form tightly bound pairs even in the absence of a Fermi surface, and these pairs behave as compact composite bosons.  
Their spatial size becomes much smaller than the average interparticle spacing, and the overlap between different pairs is negligible.  
Pair formation and condensation therefore occur as two distinct processes: pairs form at a higher temperature and subsequently undergo Bose--Einstein condensation into the lowest-energy state at a lower temperature \(T_c\).  
This behaviour can be schematically represented as:
\[
\begin{cases}
\textbf{High temperature: } T > T^{\ast} &
\Rightarrow \text{no bound pairs},\quad \Delta_0 = 0,\; n_B = 0, \\[6pt]
\textbf{Preformed pairs: } T_c < T < T^{\ast} &
\Rightarrow \text{bound pairs w/o coherence},\quad \Delta_0 > 0,\; n_0 = 0,\; \mu < 0, \\[6pt]
\textbf{Condensed phase: } T \leq T_c &
\Rightarrow \text{BEC of pairs},\quad \Delta_0 > 0,\; n_0 > 0,\; \mu < 0.
\end{cases}
\]

where \(T^{\ast}\) denotes the \textit{pairing temperature} at which bound states first appear, 
\(T_c\) is the \textit{condensation temperature}, 
\(\Delta_0\) the pair binding energy at zero momentum, 
\(n_B\) the total pair density, and \(n_0\) the condensate density.

In this regime, the centre-of-mass motion of each pair is no longer influenced by Pauli blocking, since the constituent fermions occupy bound molecular states well below the Fermi level.  
The excitation energy of a pair thus follows the free-particle, \textit{quadratic} dispersion characteristic of a boson with total mass \(2m\):
\[
    \varepsilon_K = \frac{\hbar^2 K^2}{2(2m)}.
    \label{eq:quadratic_dispersion}
\]
The chemical potential becomes negative (\(\mu < 0\)), lying below the bottom of the conduction band, as expected for a gas of bound molecules. 

Substituting \(s = 2\) and \(C_2 = \hbar^2 / (2M)\), with $M = 2m$, into the generalized expression~\eqref{eq:Tc_general_1} yields\footnote{% 
For the full derivation, see Appendix \ref{app:BCS_BEC-derivation}, section \ref{d3s2}.
} the standard result for a three-dimensional Bose gas:
\begin{equation}
    T_c = \frac{2\pi \hbar^2}{M k_B}
          \left[\frac{n_B}{\zeta(3/2)}\right]^{2/3}.
    \label{eq:Tc_strong_coupling}
\end{equation}

This quadratic dispersion corresponds to the \textit{BEC limit} of superconductivity, where tightly bound, localized pairs behave as independent bosons that condense into a single macroscopic quantum state.  



\bigskip
\bigskip
\noindent
In the previous chapters, we have developed a microscopic description of superconductivity: we explained how Cooper pairs form, how superconductivity can be viewed as a Bose--Einstein condensation of pairs, and how the system spans the weak (BCS) and strong (BEC) coupling limits.  \\
We now turn to the \emph{macroscopic} consequences of this condensate, \textit{i.e.} how a complex order parameter varies in space, how it couples to electromagnetic fields, and how measurable responses such as currents, magnetization, and characteristic length scales emerge.  
To explore these large-scale behaviours, we introduce the Landau--Ginzburg theory, a phenomenological approach that links the quantum condensate to the observable properties of real superconductors.








